\chapter{Matching Vector Codes}

\section{Introduction}

In the RM codes seen in Chapter~\ref{ldc}, we could do local correction upto $\delta<1/4$ fraction of error tolerance. Our main goal of this chapter would be to construct a code that can tolerate upto $\delta<1/2$ fraction of errors, which would be optimal for unique local decoding.

We first define what \emph{matching vectors} are, then we can specify the encoding of \emph{matching vector code}.

\begin{definition}[Matching Vectors]
    Let $S\subseteq\bbz_{m}^n$, we say that the families $\calu = \{\bfu_1, \ldots, \bfu_k\}$ and $\calv = \{\bfv_1,\ldots, \bfv_k\}$ of vectors in $\bbz_m^n$ forms an $S$-\emph{matching family} if 
    \begin{enumerate}
        \item For all $i\in[k], (\bfu_i, \bfv_i) = 0$, 
        \item For all $i\neq j$ in $[k]$, $(\bfu_j, \bfu_i) \in S$.
    \end{enumerate}
\end{definition}

\begin{greenbox}
    \begin{observation}
        Let $g$ be a generator of the unique subgroup $C_m$ of $\bbf_q^*$ of size $m$. Then for any $\bfu,\bfv,\bfw\in\bbz_m^n$ and $\lambda\in\bbz_m$, 
        \begin{align*}
            \mon_{\bfu}(g^{w+\lambda v}) 
            & = \prod_{l=1}^n \left( g^{\bfw(l) + \lambda \bfv(l)} \right)^{\bfu(l)} \\
            & = \prod_{l=1}^n g^{\bfw(l)\bfu(l)}\cdot\prod_{l=1}^n(g^{\lambda})^{\bfv(l)\bfu(l)} \\
            & = g^{(\bfu, \bfw)}\cdot (g^{\lambda})^{(\bfv, \bfu)}.
        \end{align*}
        So we can see that $M_{\bfw, \bfv}$-evaluation of $\mon_{\bfu}$ is a $C_m$-evaluation of a univariate polynomial $g^{(\bfu, \bfw)}\cdot y^{(\bfv, \bfu)}\in\bbf_q[y]$.
    \end{observation}
\end{greenbox}

Now we define the encoding of matching vector codes, but we first need to assume that there exists an $S$-matching family $\calu$ and $\calv$ of size $k$. We shall deal with this existence later. There is also an abstract decoding method mentioned that will be the main motivation behind all the decoding procedure in the coming sections.

\begin{purplebox}
   \paragraph{\large (\textsc{Encoding})} We take a $k$-length message $\bfx=(\bfx(1), \ldots, \bfx(k))\in\bbf_q^k$ and encode ot by $C_m^n$-evaluations of the polynomial 
   \begin{equation}  \label{enc}
   F_{\bfx}(z_1,\ldots,z_n) = \sum_{j=1}^{k} x(j)\mon_{u_j}(z_1,\ldots,z_n).
\end{equation}
  

    \paragraph{\large (\textsc{Abstract Decoding})} We want to retrieve $\bfx(i)$ for a specified index $i\in[k]$. The decoding procedure is:
    \begin{enumerate}
        \item Pick $\bfw\sim\mathrm{Unif}(\bbz_m^n)$. 
        \item Compute the restriction of $F \equiv F_{\bfx}$ to $M_{\bfw, v_i}$ by some queries of the corrupted $M_{w,v_i}$-evaluation of $F$ at less than or equal to $m$ points.
    \end{enumerate}
\end{purplebox}

One might ask now that what is the importance of this restriction and how it will help in retrieving $\bfx(i)$. The following observation gives us the motivation of using matching family and this second step.

\begin{greenbox}
    \begin{observation}
        The $M_{w,v_i}$-evaluation of $F$ is equivalent to the $C_m^n$-evaluation of the polynomial $f$ defined as:

        \[
            f(y) = \sum_{j=1}^{k} x(j) g^{(\bfu_j, \bfw)}\cdot y^{(\bfv_i, \bfu_j)}
        \]

        By matching family property, we can write the polynomial as:
        \[
            f(y) = x(i)g^{(\bfu_i, \bfv_i)} + \sum_{s\in S} \left(\sum_{j : (\bfu_j, \bfv_i) = s} x(j)g^{(\bfu_j, \bfv_i)} \right)y^s.
        \]  

        So we can get $\bfx(i)$ by just computing $f(0)/g^{(\bfu_i, \bfw)}$. But of course we can not construct this polynomial, so we try to find the restriction of $F$ to $M_{\bfw, \bfv_i}$.
    \end{observation}
\end{greenbox}

\section{Some Decoding Procedures}

We shall now see some explicit decoding procedures along with making the error tolerance better and better.

\subsection{Basic Decoding on Lines}

\begin{bluebox}
    \begin{proposition}
Let $\calu,\calv$ be a family of $S$-matching vectors in $\bbz_m^n$,
with $|\calu| = |\calv| = k$ and $|S| = s$. Suppose $m \mid (q-1)$, where
$q$ is a prime power. Then there exists a $q$-ary linear code encoding
$k$-long messages to $m^n$-long codewords that is
$(s+1,\delta,(s+1)\delta)$-locally decodable for all $\delta$.
\end{proposition}
\end{bluebox}

Consider the code specified in Equation~\ref{enc}. Now do the following decoding procedure:
\begin{enumerate}
    \item Pick $\bfw\sim\mathrm{Unif}(\bbz_m^n)$ and construct the multiplicative line $M_{\bfw, \bfv_i}$.
    \item Query the $M$-evaluation of $F$ at $s+1$ many consecutive locations, say $\{g^{\bfw+\lambda \bfv_i} : \lambda\in\{0,\ldots,s\}\}$ and obtain the points $c_0,\ldots,c_s$. 
    \item Interpolate the unique sparse polynomial $h\in\bbf_q[y]$ with $\mathrm{supp}(h)\subseteq S\cup\{0\}$ such that $h(g^{\lambda}) = c_{\lambda}$ for all $\lambda$.
    \item Output $h(0)/g^{(\bfw, \bfu_i)}$.
\end{enumerate}

\begin{proof}[Proof for failure probability] This is a fairly easy probability to calculate.
    \begin{align*}
        \Pr[\text{failure of decoder}] 
        & = \Pr[F(g^{\bfw + \lambda\bfv_i}) \neq c_{\lambda}\text{ for some $\lambda$}] \\
        & \leq \sum_{\lambda} \Pr[F(g^{\bfw + \lambda\bfv_i}) \neq c_{\lambda}]\\
        & \leq \sum_{\lambda}\delta = (s+1)\delta.
    \end{align*}
    
\end{proof}


\subsection{Improvement on Line Decoding}

In this section we shall see how we can make the line decoding better to get error tolerance of $\delta<1/4$ by making the matching families \emph{bounded}. 

\begin{definition}[$b$-bounded Matching Vectors]
    $\calu, \calv$ are said to be a $b$-\emph{bounded} $S$-matching family over $\bbz_m^n$ if $S$ is $b$-bounded.
\end{definition}

We now give the main proposition of this subsection analogous to Proposition~\ref{improve}.

\begin{bluebox}
\begin{proposition}
Let $0 < \sigma < 1$. Let $\mathcal{U},\mathcal{V}$ be a $\sigma m$-bounded family of $S$-matching vectors in $\mathbb{Z}_m^n$, $|\mathcal{U}| = |\mathcal{V}| = k$. Suppose $m \mid (q-1)$, where $q$ is a prime power. Then there exists a $q$-ary linear code encoding $k$-long messages to $m^n$-long codewords that is $(m,\delta,2\delta/(1-\sigma))$-locally decodable for all $\delta$.
\end{proposition}
\end{bluebox}

Consider the same code as before and the following decoding procedure:
\begin{enumerate}
    \item Pick $\bfw\sim\mathrm{Unif}(\bbz_m^n)$ and construct the multiplicative line $M_{\bfw, \bfv_i}$.
    \item Query the $M$-evaluation of $F$ at all of the $m$ locations, and obtain the points $c_0,\ldots,c_{m-1}$. 
    \item By Berlekamp-Welch algorithm, given in \cite{GRS25}, compute the unique univariate polynomial $h(y)\in\bbf_q[y]$ of degree less that $\sigma m$ and $\mathrm{supp}(h)\subseteq S\cup\{0\}$ such that $h$ agrees with the corrupted $M$-evaluation of $F$ at all but at most $m(1-\sigma)/2$ many points.
    \item Return $h(0)/g^{(\bfu_i, \bfw)}$ if such $h$ exists, else return $0$.
\end{enumerate}

\begin{proof}[Proof for failure probability]
    Take the random variable $X$ denoting the number of corrupted query points in the $M$-evaluation of $F$. Now failure occurs if $X\geq m(1-\sigma)/2$, hence failure probability is less than or equal to:
    \begin{align*}
        \Pr\left[X \geq \frac{m(1-\sigma)}{2}\right]
        & \leq \frac{1}{m(1-\sigma)/2}\bbe[X] \quad \text{[Markov Inequality]}\\
        & = \frac{1}{m(1-\sigma)/2}\sum_{\lambda\in\bbz_m}\Pr\left[F(g^{\bfw + \lambda\bfv_i})\neq c_{\lambda}\right]\\
        & \leq \frac{2m\delta}{m(1-\sigma)} = \frac{2\delta}{1-\sigma}.
    \end{align*}
\end{proof}

\begin{greenbox}
    \begin{observation}[Error Tolerance]
        As $\sigma\downarrow0$, the limiting failure probability is $2\delta$, hence we can only tolerate up to $\delta<1/4$ fraction of errors, because for error tolerance, the failure probability must be less that $1/2$. In the next subsection, we show a decoding paradigm that can tolerate up to $1/2$ fraction of errors.
    \end{observation}
\end{greenbox}

\subsection{Decoding on Multiple Lines}

Rather than constructing one random multiplicative line, we shall construct more and using the queries at each $M$-evaluation of $F$, we shall interpolate the polynomial.

\begin{bluebox}
    \begin{proposition}
Let $0 < \sigma < 1$. Let $\mathcal{U},\mathcal{V}$ be a $\sigma m$-bounded family of $S$-matching vectors in $\mathbb{Z}_m^n$, $|\mathcal{U}| = |\mathcal{V}| = k$. Suppose $m \mid (q-1)$, where $q$ is a prime power. Then for every positive integer $l$ there exists a $q$-ary linear code encoding $k$-long messages to $m^n$-long codewords that for all $\delta < (1-\sigma)/2$ is $(lm,\delta,\exp_{\delta,\sigma}(-l))$-locally decodable.
\end{proposition}
\end{bluebox}

Consider the same code as before and fix an $l\in\bbn$. Now here we shall not construct the polynomial like we have do before, we shall consider $l$ many $M$-evaluations of $F$ and instead of interpolating the polynomial directly, we shall assign weight to each field element $e$ with respect to the likelihood of $x(i) = e$. And from that we shall return the unique highest weight field element, if exists. The decoding procedure is the following:

\begin{enumerate}
    \item Take $\bfw_1,\ldots,\bfw_l$ uniformly at random from $\bbz_m^n$ and consider the $l$ many multiplicative lines $\{M_{\bfw_j, \bfv_i} : j\in[l]\}$.
    \item Query at the $lm$ many coordinates of all the $M$-evaluations of $F$ and obtain $l$ many vectors in $\bbf_q^m$, say $\bfc^1,\ldots,\bfc^m$. For each $e\in\bbf_q$, compute:
    \[
        \mathrm{weight}(e) = \sum_{j=1}^{l}\mathrm{weight}(e,M_{\bfw_j, \bfv_i}).
    \]
    \item Return $e$ with the highest weight. If it is not unique, then return $0$.
\end{enumerate}

\begin{greenbox}
    \begin{observation}
        Suppose there are two distinct $e_1,e_2\in\bbf_q$ such that $\mathrm{weight}(e_1), \mathrm{weight}(e_2)>lm(1+\sigma)/2$. Then by PHP, there would be two distinct polynomials $h_1, h_2\in\bbf_q[y]$, both of degree less than $\sigma m$ and two multiplicative lines $M_1, M_2$ such that $\mathrm{weight}(e_1,M_1)>m(1+\sigma)/2$ and $\mathrm{weight}(e_2,M_2)>m(1+\sigma)/2$. Now by inclusion-exclusion, the numeber of agreement points of $h_1$ and $h_2$ is at least $2m(1+\sigma)/2 - m = \sigma m$. Thus we can not have such a pair.
    \end{observation}
\end{greenbox}

\begin{proof}[Proof for failure probability]
    Consider the random variable $X$ denoting the number of corrupted queried points and failure probability is less than or equal to $\Pr[X\geq lm(1+\sigma)/2]$. Note that, 
    \[
        X = \sum_{j = 1}^{l}\left( \sum_{\lambda\in\bbz_m} 
\mathbbm{1}_{\{F(g^{(\bfw_j,\bfv_i)}_\lambda)\neq \bfc^j_{\lambda}\}}
\right)
    \]

    Now the expected value of $X$ is:
    \begin{align*}
    \bbe[X]
    &= \sum_{j = 1}^{l}\left( \sum_{\lambda\in\bbz_m} 
    \Pr[F(g^{(\bfw_j,\bfv_i)}_\lambda)\neq \bfc^j_{\lambda}]
    \right)\\
    &\leq \sum_{j = 1}^{l}\left( \sum_{\lambda\in\bbz_m} \delta \right) = ml\delta.
\end{align*}

Now $X$ is the sum of $lm$ many independent and identically distributed Bernoulli random variable of success probability $\delta$. Now we can apply \emph{Chernoff bound}, given in \cite{Alon2000}, in the following way to find an upper bound for failure probability. Using $\delta<(1 - \sigma)/2$,
\begin{align*}
    \Pr[X\geq lm(1+\sigma)/2] 
    & \leq \Pr[X - \bbe[X] \geq lm(1+\sigma)/2 - lm(1 - \sigma)/2] \\
    &  = \Pr[X - \bbe[X]\geq ml\sigma] \\
    & \leq \exp_{\sigma, \delta}(-l) \quad \text{Chernoff Bound}. 
\end{align*}
\end{proof}

Now we can show how we increase the error tolerance up to $\delta < 1/2$. Since the proposition holds for every $\delta < (1-\sigma)/2$, letting $\sigma \downarrow 0$ yields $\delta \uparrow 1/2$. Hence the decoder tolerates any error fraction $\delta < 1/2$.  


